% 11. Скворцова 10.14357/08696527200411



 \def\stat{skvortsova}



\def\tit{ПРОГРАММНЫЙ КОМПЛЕКС ДЛЯ~ПОЗИЦИОНИРОВАНИЯ 

АКУСТИЧЕСКИХ ДОННЫХ СИСТЕМ}



\def\titkol{Программный комплекс для~позиционирования 

акустических донных систем}





\def\aut{В.\,А.~Смирнов$^1$, Н.\,Н.~Скворцова$^2$, Е.\,М.~Кончеков$^3$, 

В.\,А.~Ларичев$^4$, Г.\,А.~Максимов$^5$}



\def\autkol{В.\,А.~Смирнов, Н.\,Н.~Скворцова, Е.\,М.~Кончеков и~др.}

%, В.\,А.~Ларичев$^4$, Г.\,А.~Максимов$^5$}



\titel{\tit}{\aut}{\autkol}{\titkol}



\index{Смирнов В.\,А.}

\index{Скворцова Н.\,Н.}

\index{Кончеков Е.\,М.}

\index{Ларичев В.\,А.}

\index{Максимов Г.\,А.}

\index{Smirnov V.\,A.}

\index{Skvortsova N.\,N.}

\index{Konchekov E.\,M.}

\index{Larichev V.\,A.}

\index{Maximov G.\,A.}





%{\renewcommand{\thefootnote}{\fnsymbol{footnote}}

%\footnotetext[1]

%{Статья подготовлена при частичной поддержке РФФИ (проекты 18-29-03124-мк 

%и~18-29-03081-мк).}} 





\renewcommand{\thefootnote}{\arabic{footnote}}

\footnotetext[1]{Акустический институт им.\ академика Н.\,Н.~Андреева; 

МИРЭА~--- Российский технологический университет, \mbox{Viperx15@mail.ru}}

\footnotetext[2]{Институт общей физики им.\ А.\,М.~Прохорова Российской 

академии наук; МИРЭА~--- Российский технологический университет; Национальный 

исследовательский ядерный университет <<МИФИ>>, \mbox{mukudori@mail.ru}}

\footnotetext[3]{Институт общей физики им.\ А.\,М.~Прохорова Российской 

академии наук; Российский национальный исследовательский медицинский университет 

им.\ Н.\,И.~Пирогова, \mbox{eukmek@gmail.com}}

\footnotetext[4]{Акустический институт им.\ академика Н.\,Н.~Андреева; 

\mbox{larido@gmail.com}}

\footnotetext[5]{Акустический институт им.\ академика Н.\,Н.~Андреева; 

\mbox{gamaximov@mail.ru}}



 \label{st\stat}



 \thispagestyle{headings}

 

 \vspace*{-12pt} 



     \Abst{На основе алгоритмов корреляционного и~спектрального анализа разработан 

программный комплекс для обработки радиофизических сигналов с~целью точного 

позиционирования приемных элементов донных сейсмических кос. Сейсмокосы были 

испытаны в~Геленджикской бухте на Черном море. Представлен результат 

позиционирования системы, полученный с~помощью разработанных алгоритмов по 

измеренным акустическим дистанциям. Произведено сравнение полученных результатов 

с~контрольными данными измерений эхолота. Показано, что измерение акустических 

дистанций при использовании разработанного программного комплекса отвечает 

требованиям для позиционирования систем морской 2D-сейсморазведки.}

     

     \KW{гидроакустика; мониторинг; морская сейсморазведка; твердотельная донная 

цифровая сейсмокоса; позиционирование; корреляционный анализ; спектральный анализ}

     

\DOI{10.14357\!/\!08696527200411} 



\vskip 10pt plus 9pt minus 6pt



 \thispagestyle{headings}

 

 \vspace*{-18pt}

     

\section{Введение}



    В настоящее время одной из наиболее актуальных задач гидроакустики 

стала задача сейсмической разведки для целей обнаружения полезных 

ископаемых на шельфах. Морская сейсморазведка обычно проводится 

следующими способами: раскладкой донных сейсмических кос, 

заглублением множества отдельных маяков-ответчиков, буксировкой 

сейсмических кос у поверхности (глубина $\sim 5\mbox{--}10$~м)~[1]. 

Получение сейсмических данных при буксировке сейсмических кос у 

поверхности моря является актуальной как с~научной, так и~с~практической 

точки зрения задачей~[2]. Решение данной задачи связано с~разработкой 

сложных программных комплексов для измерения, обработки и~хранения 

экспериментальных данных~[3]. Методы детектирования и~обработки таких 

радиофизических сигналов необходимы для исследования сигналов 

и~в~других областях науки и~техники: в~акустике, геофизике, исследовании 

плазменной~[4--7] и~атмосферной турбулентности~[8] и~т.\,д.

    

    Современной и~актуальной стала задача обработки лоцирующих 

акустических сигналов донных и~буксируемых акустических антенн 

(сейсмокос) для получения массива расстояний между датчиками 

сейсмокосы и~источником лоцирующих сигналов для позиционирования 

системы~[9]. В~статье будет рассмотрен алгоритм работы созданного 

программного обеспечения для обработки записей ло\-ци\-ру\-ющих сигналов на 

приемниках сейсмокосы, выделения времен распространения акустического 

сигнала по прямому лучу от источника с~известными координатами до 

приемного элемента сейсмокосы для измерения акустических дистанций, 

позиционирования и~вычисления ориентации приемных элементов 

в~пространстве.



\section{Позиционирование донной сейсмокосы}



В Акустическом институте им.\ академика Н.\,Н.~Андреева (АО АКИН) 

разработаны и~производятся цифровые твердотельные донные сейсмокосы 

для инженерной морской сейсморазведки и~мониторинга. После заглубления 

такой косы перед проведением сейсмических измерений требуется 

определить расположение и~ориентацию каждого приемного элемента. Для 

решения данной задачи в~районе раскладки и~после заглубления 

сейсмокосы используется специальное суд\-но-пин\-гер\footnote{Пингер~--- 

это судно, осуществляющее проход по заданной траектории раскладки погружаемых 

акустических средств (в~данном случае сейсмокосы) и~излучающее широкополосные 

мощные лоцирующие импульсы для позиционирования данных средств.}, оснащенное 

спаркером, предназначенным для сейсмических измерений. Для излучения 

локационных импульсов судно оборудовано высокочастотным 

пьезокерамическим излучателем, работающим в~полосе частот 

$\sim30$~кГц, с~помощью которого производится излучение серии 

высокочастотных широкополосных импульсов. Запись принятых сигналов 

с~приемных элементов донной сейсмокосы синхронизирована с~началом 

излучения  

суд\-на-пин\-гера. 

     

     Работа программы начинается с~выбора пользователем директории, 

содержащей записи сигналов приемных элементов сейсмокосы в~формате 

SEGY~2.0 и~записи координат GPS траектории движения суд\-на-пин\-ге\-ра 

с~точками, соответствующими моментам излучения лоцирующего сигнала 

(хлопка). Общая траектория движения суд\-на-пин\-ге\-ра и~моменты 

времени, соответствующие записям с~сейсмокосы, приведены на рис.~1. На 

рис.~1 символом~<<*>> отмечен момент излучения лоцирующего импульса. 

Значком~<<$\square$>> отмечен момент излучения лоцирующего сигнала, 

для которого в~заданной директории была найдена запись принятого сигнала 

сейсмокосы. Каждому значку~<<$\square$>> соответствует участок записи 

временн$\acute{\mbox{ы}}$х выборок всех приемных элементов сейсмокосы в~одном из 

файлов SEGY. 



\begin{figure} %fig1

\vspace*{1pt}

\begin{center}

 \mbox{%

 \epsfxsize=100mm 

\epsfbox{skv-1.eps}

 }

 \end{center}

\vspace*{-6pt}

\Caption{Траектория движения судна-пингера. По осям~$X$ и~$Y$ отмечено расстояние рабочего 

региона в~метрах. Общее число точек траектории~--- 6444. Число моментов излучения 

лоцирующего импульса~--- 1244 (значок~<<*>>). Число моментов излучения, для которых 

проводится анализ записей сигналов~--- 369 (значок~<<$\square$>>)}

\vspace*{-2pt}

\end{figure}



     Один приемный элемент сейсмокосы имеет~5~приемных каналов 

(3~низ\-ко\-час\-тот\-ных для сейсмических измерений и~2~высокочастотных для 

позиционирования и~поиска ориентации). Частота дискретизации для всех 

приемных каналов соответствует 2~кГц. На рис.~2 и~3 представлены 

временн$\acute{\mbox{ы}}$е выборки записей всех высокочастотных приемных элементов 

сейсмокосы при излучении одной серии импульсов, загруженные 

программой из файлов SEGY. Построение записей сигналов проводится 

вертикально для визуализации максимального числа каналов на экране 

пользователя. На рис.~3 представлен небольшой участок для рассмотрения 

формы лоцирующих импульсов и~результатов детектирования. Каждому 

приемному элементу (гидрофону) соответствуют два приемных канала, что 

необходимо для дальнейшего вычисления его ориентации в~пространстве.

     

\begin{figure} %fig2

\vspace*{1pt}

\begin{center}

 \mbox{%

 \epsfxsize=100mm 

\epsfbox{skv-2.eps}

 }

 \end{center}

\vspace*{-6pt}

\Caption{Запись сигналов приемных элементов сейсмокосы при излучении одной серии 

импульсов (приближен участок записи  

каналов 138--168, только четные каналы)}

%\end{figure}

\vspace*{8pt}

%\begin{figure} %fig3

\vspace*{1pt}

\begin{center}

 \mbox{%

 \epsfxsize=100mm 

\epsfbox{skv-3.eps}

 }

 \end{center}

\vspace*{-6pt}

\Caption{Финальный результат позиционирования приемных элементов сейсмокосы. По 

осям $X$ и~$Y$ отмечено расстояние рабочего региона в~метрах. Черным цветом 

представлены вычисленные координаты приемных каналов 352--510. Серым цветом 

представлены измеренные (значок~<<$\square$>>) и~скорректированные в~пределах 

погрешности координаты GPS-маяка в~моменты излучения лоцирующего импульса 

(значок~<<$\times$>>)}

\vspace*{-2pt}

\end{figure}



    На рис.~2 используются следующие обозначения:

     \begin{itemize}

\item точками на временн$\acute{\mbox{ы}}$х выборках показаны рассчитанные по данной 

записи времена распространения лоцирующего сигнала по прямому лучу. 

На рис.~2 приведен пример некорректного детектирования (отброшенное 

измерение), не прошедшего валидацию (хотя бы по одному критерию). 

Дистанции, полученные из таких времен, не принимают участия в~расчете. 

Изменяя пороговые значения отсева результатов, можно увеличить число 

измеренных дистанций, но с~уменьшением пороговых значений 

повышается вероятность попадания некорректно измеренной дистанции 

в~расчет позиционирования; 

\item значком <<$\times$>>отмечены времена распространения, восстановленные из уже 

рассчитанного расположения элементов (по всем $\sim1200$~записям), они появляются после 

первого предварительного анализа всех записей и~вычисления предварительного результата 

позиционирования;

\item символом <<$+$>> показаны начало и~конец корреляционных 

интервалов и~значения корреляционных коэффициентов сигналов двух 

приемных каналов каждого элемента, необходимые для поиска ориентации 

приемного элемента в~пространстве, цифрой указан уровень корреляции 

сигналов.

\end{itemize}



     

     Расчет времен прямого распространения лоцирующего сигнала 

проводится следующими методами численного анализа временн$\acute{\mbox{ы}}$х выборок.

     \begin{enumerate}[1.]

  \item Вычисление отношения сигнал\!/\!шум на всем интервале сигнала для 

выбора рабочего участка.

  \item Корреляция огибающей принятого сигнала и~огибающей излученной 

последовательности лоцирующих импульсов. Значения ниже порогового 

уровня корреляции отбрасываются. 

\item Медианная фильтрация\footnote{Медианный фильтр~--- один из видов цифровых 

фильтров, широко используемый в~цифровой обработке сигналов и~изображений для 

уменьшения уровня шума. Медианный фильтр является нелинейным фильт\-ром c конечной импульсной характеристикой. 

Значения отсчетов внутри окна фильтра сортируются в~порядке возрастания (убывания); 

и~значение, находящееся в~середине упорядоченного списка, поступает на выход 

фильтра. В~случае четного числа отсчетов в~окне выходное значение фильтра равно 

среднему значению двух отсчетов в~середине упорядоченного списка. Окно 

перемещается вдоль фильтруемого сигнала, и~вычисления повторяются.} принятого 

сигнала с~заданным окном для коррекции времен распространения сигналов и~борьбы 

с~замираниями.

     \end{enumerate}

     

     Решение задачи позиционирования проходит в~два этапа.

  \begin{enumerate}[1.]

  \item Вычисление предварительного результата с~помощью метода 

Ньютона или трилатерации.

\item Вычисление конечного результата с~помощью итерационных методов: градиентного 

спуска\footnote{Градиентный спуск~--- метод нахождения локального экстремума функции с помощью движения вдоль градиента.} 

или метода сопряженных градиентов\footnote{Метод сопряженных градиентов~--- метод 

нахождения локального экстремума функции на основе информации о ее значениях и~ее градиенте.}.

     \end{enumerate}

     

     После запуска расчета программа автоматически анализирует все 

со\-по\-став\-лен\-ные моменты излучения с~записями сигналов, проводит 

измерение акус\-ти\-че\-ских дистанций и~представляет результат 

позиционирования. На рис.~3 представлен конечный результат 

позиционирования приемных элементов донной сейсмокосы. 



\section{Ориентация приемных элементов донной 

сейсмокосы}



     Задача поиска ориентации решается путем построения математической 

модели распространения акустического сигнала с~известными частотными 

характеристиками от источника к~двум гидрофонам приемника, 

находящимся на известном расстоянии друг от друга (14--17~см). 

     Затем осуществляется перебор значений входных параметров модели:

     \begin{itemize}

\item горизонтальное расположение приемника (перебор с~точностью 

0,5$^\circ$);

\item вертикальное расположение приемника (перебор с~точностью 

0,5$^\circ$);

\item дистанция между приемниками одного элемента (перебор 

с~точностью 0,5~мм);

\item частота несущей излучаемого сигнала ($\sim30$~кГц).

     \end{itemize}

     

     Цель перебора для данной задачи заключается в~поиске параметров 

модельной функции, для которой корреляционные коэффициенты сигналов 

приемных каналов одного элемента из разных источников совпадают 

в~пределах допустимой погрешности с~экспериментальными значениями, 

определяемыми в~ходе анализа сигналов приемных каналов после решения 

задачи позиционирования.

     

     Расположение элемента, при котором модельные коэффициенты 

совпадают с~экспериментальными в~пределах допустимой погрешности, 

считается правильным ответом. Если такое расположение не было найдено 

при переборе всех доступных параметров моделирования, ориентация 

остается не определенной и~отмечается серым цветом.

     

     Результат сравнения полученной математической модели, 

учитывающей расположение приемного элемента, с~исходными 

(экспериментальными) коэффициентами представлен на рис.~4. 

Коррелируемые сигналы представлены справа.

     

\begin{figure}[b] %fig4

\vspace*{1pt}

\begin{center}

 \mbox{%

 \epsfxsize=100mm 

\epsfbox{skv-4.eps}

 }

 \end{center}

\vspace*{-6pt}

\Caption{Сравнение функции распределения взаимно корреляционных коэффициентов приемных 

каналов одного элемента (модель и~эксперимент)}

\end{figure}

     

     На рис.~5 представлен результат вычисления ориентации приемных 

элементов сейсмокосы в~пространстве (серым цветом). Серое значение 

соответствует дистанции между приемными каналами приемника в~метрах 

(на этапе моделирования это один из параметров, который может изменяться 

в~небольших пределах).

     

     Таким образом, можно применить разработанные алгоритмы 

спектрально-кор\-ре\-ля\-ци\-он\-ной обработки радиофизических сигналов для 

задачи детектирования времен прямого распространения акустических 

лоцирующих сигналов для решения задачи позиционирования и~ориентации 

приемных элементов донной сейсмокосы. 

     

     Полученные результаты позиционирования сравнивались с~данными 

глубин эхолота, установленного на суд\-не-пин\-ге\-ре. Конечная ошибка 

позиционирования не может быть меньше точности GPS-ма\-яка, 

установленного на судне (1--3~м), однако при большом числе точных 

измерений дистанций (более~50 для каждого приемного элемента) 

статистически ошибка позиционирования может быть уменьшена до~0,5~м 

и~менее. При расчете позиционирования существует возможность 

зафиксировать глубину, замещая ее данными эхолота, что может уменьшить 

ошибку позиционирования до~0,1~м.



\begin{figure} %fig5

\vspace*{1pt}

\begin{center}

 \mbox{%

 \epsfxsize=100mm 

\epsfbox{skv-5.eps}

 }

 \end{center}

\vspace*{-6pt}

\Caption{Ориентация элементов донной сейсмокосы}

\end{figure}



\section{Заключение}



    Реализовано бортовое кроссплатформенное бортовое программное 

обеспечение (протестировано на Windows XP, Windows~7, Windows~10, 

Ubuntu~18, Debian~9 и~10) для визуализации и~сопоставления данных 

траекторий движения суд\-на-пин\-ге\-ра, обработки и~анализа данных 

акустических сигналов донных сейсмокос, позиционирования и~поиска 

ориентации приемных элементов донных сейсмокос. Программное 

обеспечение создавалось средствами С++~17 с~по\-мощью открытого пакета 

для разработки QT Creator~5. Реализованы алгоритмы для выделения времен 

прямого распространения акустических сигналов от источника с~известными 

координатами (пингера) до приемных элементов донной сейсмокосы, 

реализован ряд алгоритмов для борьбы с~замираниями и~помехами. 

Проведено сравнение полученных результатов позиционирования 

с~данными эхолота. Полученная точность позиционирования достаточна для 

задач морской 2D-сейс\-мо\-раз\-ведки.



{\small\frenchspacing

{%\baselineskip=10.8pt

\addcontentsline{toc}{section}{Литература}

\begin{thebibliography}{9}



\bibitem{1-skv}

\Au{Максимов Г.\,А., Денисов Д.\,Н., Ларичев В.\,А., Лесонен~Д.\,Н., Григорьев~А.\,Г., 

Корольков~З.\,А.} Цифровая твердотельная буксируемая сейсмокоса малого диаметра для 

морской инженерной сейсморазведки~/\!\!/ Учен. зап. физ. фак-та Моск. ун-та, 2017. №\,5. 

Ст. №\,1750126.

\bibitem{2-skv}

\Au{Ампилов Ю.\,П.} Освоение шельфа Арктики и~Дальнего Востока. Проблемы 

и~перспективы~/\!\!/ Offshore Russia, 2014. №\,4. С.~8--15.

\bibitem{3-skv}

\Au{Малахов Д.\,В., Смирнов В.\,А., Скворцова Н.\,Н., Соколов~А.\,С., Степахин~В.\,Д.,

Борзосеков~В.\,Д., Харчев~Н.\,К.} Радиофизический про\-грам\-мно-ап\-па\-рат\-ный комплекс сбора и~обработки данных на 

стеллараторе Л-2М~/\!\!/ Инженерная физика, 2015. №\,8. С.~38--44.



\bibitem{5-skv} %4

\Au{Skvortsova N.\,N., Akulina D.\,K., Batanov G.\,M., Kharchev~N.\,K., Kolik~L.\,V., 

Kovrizhnykh~L.\,M., Letunov~A.\,A., Logvinenko~V.\,P., Malakhov~D.\,V., Petrov~A.\,E., 

Pshenichnikov~A.\,A., Sarksyan~K.\,A., Voronov~G.\,S.} Effect of ECRH regime on 

characteristics of short-wave turbulence in plasma of the L-2M~/\!\!/ Plasma Phys. 

Contr.~F., 2010. Vol.~52. Art. No.\,055008. 

\bibitem{6-skv} %5

\Au{Батанов Г.\,М., Борзосеков В.\,Д., Колик Л.\,В., Малахов~Д.\,В., Петров~А.\,Е., 

Пшеничников~А.\,А., Сарксян~К.\,А., Скворцова~Н.\,Н., Харчев~Н.\,К.} Длинноволновая 

турбулентность в~плазме стелларатора Л-2М при элект\-рон\-но-цик\-ло\-трон\-ном 

нагреве~/\!\!/ Вопросы атомной науки и~техники. Сер.: Термоядерный синтез, 2011. №\,2. 

С.~70--75. 

\bibitem{4-skv} %6

\Au{Батанов Г.\,М., Борзосеков В.\,Д., Коврижных Л.\,М., Колик~Л.\,В., Кончеков~Е.\,М., 

Малахов~Д.\,В., Петров~А.\,Е., Сарксян~К.\,А., Скворцова~Н.\,Н., Степахин~В.\,Д., 

Харчев~Н.\,К.} Рассеяние назад излучения гиротрона при ЭЦ нагреве плазмы на 

стеллараторе Л-2М и~коротковолновая турбулентность~/\!\!/ Физика плазмы, 2013. Т.~39. 

№\,6. С.~511--522.



\bibitem{7-skv}

\Au{Скворцова Н.\,Н., Степахин В.\,Д., Малахов Д.\,В., Сорокин~А.\,А., Батанов~Г.\,М., 

Борзосеков~В.\,Д., Глявин~М.\,Ю., Колик~Л.\,В., Кончеков~Е.\,М., Летунов~А.\,А., 

Петров~А.\,Е., Рябикина~И.\,Г., Сарксян~К.\,А., Соколов~А.\,С., Смирнов~В.\,А., 

Харчев~Н.\,К.} Создание рельефа на молибденовых пластинах в~разрядах, ини\-ци\-иру\-емых 

излучением гиротрона в~порошках ме\-талл--ди\-элект\-рик~/\!\!/ Известия вузов. 

Радиофизика, 2015. Т.~58. №\,9. C.~779--788.

\bibitem{8-skv}

\Au{Маслов С.\,А., Харчевский А.\,А., Смирнов В.\,А.} Применение вейвлета Хаара для 

анализа плазменных и~атмосферных флуктуаций~/\!\!/ Ядерная физика и~инжиниринг, 

2016. Т.~7. №\,5. C.~448--452.

\bibitem{9-skv}

\Au{Смирнов В.\,А., Скворцова Н.\,Н., Максимов~Г.\,А., Ларичев~В.\,А., Смагин~Д.\,А., 

Лекомцев~В.\,М.} Алгоритмы обработки радиофизических сигналов для систем 

дистанционного мониторинга в~гидросфере~/\!\!/ Прикладная физика, 2019. №\,5. 

C.~85--92.

\end{thebibliography}

} }





\vspace*{-3pt}



\hfill{\small\textit{Поступила в~редакцию 20.02.20}}









\vspace*{6pt}



\hrule



\vspace*{2pt}



%\newpage



%\vspace*{-28pt}



\hrule



\vspace*{8pt}



\def\tit{SOFTWARE PACKAGE FOR~POSITIONING\\ OF~ACOUSTIC BOTTOM SYSTEMS}



\def\titkol{Software package for positioning of acoustic bottom systems}



\def\aut{V.\,A.~Smirnov$^{1,2}$, N.\,N.~Skvortsova$^{2,3,4}$, 

E.\,M.~Konchekov$^{3,5}$, V.\,A.~Larichev$^1$, and~G.\,A.~Maximov$^1$}



\def\autkol{V.\,A.~Smirnov, N.\,N.~Skvortsova, E.\,M.~Konchekov, et al.}

%, V.\,A.~Larichev$^1$, and~G.\,A.~Maximov}



\titel{\tit}{\aut}{\autkol}{\titkol}



\vspace*{-11pt}





\noindent

$^1$N.\,N.~Andreev Acoustic Institute, 4~Shvernika Str., Moscow 117036, Russian\linebreak

$\hphantom{^1}$Federation







\noindent

$^2$MIREA~--- Russian Technological University, 78 Vernadskogo Ave., Moscow\linebreak

$\hphantom{^1}$119454, Russian 

Federation



\noindent

$^3$Prokhorov General Physics Institute of the Russian Academy of Sciences, 38~Va-\linebreak

$\hphantom{^1}$vilov Str., Moscow  119991, Russian Federation



\noindent

$^4$National Research Nuclear University ``MEPhI,'' 31~Kashirskoe Shosse,\linebreak

$\hphantom{^1}$Moscow 115409, 

Russian Federation



\noindent

$^5$Pirogov Russian National Research Medical University, 1~Ostrovitianov Str.,\linebreak

$\hphantom{^1}$Moscow 117997,  Russian Federation





\def\leftfootline{\small{\textbf{\thepage}

\hfill Sistemy i Sredstva Informatiki~---

Systems and Means of Informatics 2020 vol~30 no\,4}

}%

 \def\rightfootline{\small{Sistemy i Sredstva Informatiki~---

 Systems and Means of Informatics 2020 vol~30 no\,4

\hfill \textbf{\thepage}}}

 

 

\vspace*{3pt}



\Abste{The problem of accurate positioning of the receiving elements of the bottom seismic streamers is 

considered. A~new software package for processing, based on the algorithms of correlation and spectral 

analysis, of radio-physical signals is developed. Seismic streamers were tested in the Gelendzhik bay on 

the Black Sea. The system positioning results, obtained using the developed algorithms according to the 

measured acoustic distances, are presented and compared with the echo sounder measurement data. It is 

shown that measurements of acoustic distances using the developed software meet the requirements for 

positioning of two-dimensional marine seismic systems.}



\KWE{hydroacoustics; monitoring; marine seismic; solid-state bottom digital seismic streamer; 

positioning; correlation analysis; spectral analysis}



\DOI{10.14357\!/\!08696527200411}



%\vspace*{-24pt}



%\Ack

%\noindent

 



\vspace*{-12pt}



\renewcommand{\bibname}{\protect\rmfamily References}

%\renewcommand{\bibname}{\large\protect\rm References}



{\small\frenchspacing

{%\baselineskip=10.8pt

\addcontentsline{toc}{section}{References}

\begin{thebibliography}{9}



\bibitem{1-skv-1}

\Aue{Maximov, G.\,A., D.\,M. Denisov, V.\,A. Larichev, D.\,N.~Lesonen, A.\,G.~Grigor'ev, and 

Z.\,A.~Korolkov.} 2017. Tsifrovaya tverdotel'naya buksiruyemaya seysmokosa malogo diametra dlya 

morskoy inzhenernoy seysmorazvedki [Digital solid seismic streamer of a~small diameter for the marine 

engineering seismic exploration]. \textit{Uchen. zap. fiz. fak-ta Mosk. un-ta} [Moscow University Physics 

Bull.] 5:1750126.

\bibitem{2-skv-1}

\Aue{Ampilov, Yu.\,P.} 2014. Osvoenie shel'fa Arktiki i~Dal'nego Vostoka. Problemy i~perspektivy 

[Developing the arctic shelf and far eastern shelf. Problems and perspectives]. \textit{\\Offshore Russia} 

4:8--15.

\bibitem{3-skv-1}

\Aue{Malakhov, D.\,V., V.\,A. Smirnov, N.\,N. Skvortsova, A.\,S.~Sokolov, V.\,D.~Stepakhin, 

V.\,D.~Borzosekov, and N.\,K.~Harchev.} 2015. Radiofizicheskiy apparatnyy kompleks sbora i~obrabotki dannykh na 

stellaratore L-2M [Radiophysical hardware and software system for acquisition and

data processing on  stellarator L-2M]. 

\textit{Inzhenernaya fizika} [Engineering Physics] 8:38--44.



\bibitem{5-skv-1} %4

\Aue{Skvortsova, N.\,N., D.\,K.Akulina, G.\,M.~Batanov, N.\,K.~Kharchev, L.\,V.~Kolik, 

L.\,M.~Kovrizhnykh, A.\,A.~Letunov, V.\,P.~Logvinenko, D.\,V.~Malakhov, A.\,E.~Petrov, 

A.\,A.~Pshenichnikov, K.\,A.~Sarksyan, and G.\,S.~Voronov}. 2010. Effect of ECRH regime on 

characteristics of short-wave turbulence in plasma of the L-2M stellarator. \textit{Plasma Phys. 

Control.~F.} 52:055008.

\bibitem{6-skv-1} %5

\Aue{Batanov, G.\,M., V.\,D. Borzosekov, L.\,V.~Kolik, D.\,V.~Malakhov, A.\,E.~Petrov, 

A.\,A.~Pshenichnikov, K.\,A.~Sarksyan, N.\,N.~Skvortsova, and N.\,K.~Kharchev.} 2011. 

Dlinnovolnovaya turbulentnost' v~plazme stellaratora L-2M pri elektronno-tsiklotronnom nagreve 

[Long-wavelength turbulence during ECR heating of the plasmain L-2M stellarator]. \textit{Voprosy 

atomnoy nauki i~tekhniki. Ser. Termoyadernyy sintez} [Problems of Atomic Science and Technology. 

Thermonuclear Fusion Ser.] 2:70--75.

\bibitem{4-skv-1} %6

\Aue{Batanov, G.\,M., V.\,D. Borzosekov, L.\,M. Kovrizhnykh, L.\,V.~Kolik, E.\,M.~Konchekov, 

D.\,V.~Malakhov, A.\,E.~Petrov, K.\,A.~Sarksyan, N.\,N.~Skvortsova, V.\,D.~Stepakhin, and 

N.\,K.~Kharchev}. 2013. Backscattering of gyrotron radiation and short-wavelength 

turbulence during electron cyclotron resonance plasma heating in the L-2M stellarator. 

\textit{Plasma Phys. Rep.} 39(6):444--455.

\bibitem{7-skv-1}

\Aue{Skvortsova, N.\,N., V.\,D. Stepakhin, D.\,V.~Malakhov, A.\,A.~Sorokin, G.\,M.~Batanov, 

V.\,D.~Borzosekov, M.\,Yu.~Glyavin, L.\,V.~Kolik, E.\,M.~Konchekov, A.\,A.~Letunov, A.\,E.~Petrov, 

I.\,G.~Ryabikina, K.\,A.~Sarksyan, A.\,S.~Sokolov, V.\,A.~Smirnov, and N.\,K.~Kharchev}. 2015. 

Relief creation on molybdenum plates in discharges initiated by gyrotron 

radiation in metal--dielectric powder mixtures. \textit{Radiophys. 

Quantum El.} 58(9):701--709.

\bibitem{8-skv-1}

\Aue{Maslov, S.\,A., A.\,A. Kharchevsky, and V.\,A.~Smirnov.} 2017. 

Application of the Haar wavelet to the analysis of 

plasma and atmospheric fluctuations. \textit{Phys. Atom. Nucl.} 

80(11): 1692--1696.

\bibitem{9-skv-1}

\Aue{Smirnov, V.\,A., N.\,N. Skvortsova, G.\,A. Maksimov, V.\,A.~Larichev, D.\,A.~Smagin, and 

V.\,M.~Lekomcev.} 2019. Algoritmy obrabotki radiofizicheskikh signalov dlya sistem distantsionnogo 

monitoringa v~gidrosfere [Data processing algorithms of radiophysical signals for remote monitoring 

systems in hydrosphere]. \textit{Prikladnaya fizika} [Applied Physics] 5:85--92.

\end{thebibliography}

} }



%\vspace*{-3pt}



\hfill{\small\textit{Received February 20, 2020}} 



%\vspace*{-16pt}





\Contr



\noindent

\textbf{Smirnov Vitaliy A.} (b.\ 1990)~--- PhD student, Department of Modeling Radio-Physical 

Processes, MIREA~--- Russian Technological University, 78~Vernadskogo Ave., Moscow 119454, 

Russian Federation; leading engineer, Department of Acoustic Transducers, N.\,N.~Andreev Acoustic 

Institute, 4~Shvernika Str., Moscow 117036, Russian Federation; \mbox{Viperx15@mail.ru}



\vspace*{3pt}



\noindent

\textbf{Skvortsova Nina N.} (b.\ 1956)~--- Doctor of Science in physics and mathematics, professor, 

Department of Modeling Radio-Physical Processes, MIREA~--- Russian Technological University, 

78~Vernadskogo Ave., Moscow 119454, Russian Federation; professor, National Research Nuclear 

University ``MEPhI,'' 31~Kashirskoe Shosse, Moscow 115409, Russian Federation; leading scientist, 

Prokhorov General Physics Institute of the Russian Academy of Sciences, 38~Vavilov Str., Moscow 

119991, Russian Federation; \mbox{mukudori@mail.ru}



\vspace*{3pt}



\noindent

\textbf{Konchekov Evgeny M.} (b.\ 1988)~--- Candidate of Science (PhD) in physics and mathematics, 

associate professor, Department of Physics, Pirogov Russian National Research Medical University, 

1~Ostrovitianov Str., Moscow 117997, Russian Federation; senior scientist, Prokhorov General Physics 

Institute of the Russian Academy of Sciences, 38~Vavilov Str., Moscow 119991, Russian Federation; 

\mbox{eukmek@gmail.com}



\vspace*{3pt}



\noindent

\textbf{Larichev Vladimir A.} (b.\ 1972)~--- Candidate of Science (PhD) in physics and mathematics, 

scientist, N.\,N.~Andreev Acoustic Institute, 4~Shvernika Str., Moscow 117036, Russian Federation; 

\mbox{larido@gmail.com}



\vspace*{3pt}



\noindent

\textbf{Maksimov German A.} (b.\ 1963)~--- Doctor of Science in physics and mathematics, head of  

department, N.\,N.~Andreev Acoustic Institute, 4~Shvernika Str., Moscow 117036, Russian Federation; 

\mbox{gamaximov@mail.ru}



\label{end\stat}



\renewcommand{\bibname}{\protect\rm Литература} 

